\documentclass[12pt,a4paper]{article}

\usepackage[margin=2.5cm]{geometry}
\usepackage{amsmath}
\usepackage{amssymb}
\usepackage{graphicx}
\usepackage{booktabs}
\usepackage{hyperref}
\usepackage{parskip}
\usepackage{titlesec}
\usepackage{xcolor}
\usepackage{mdframed}

\hypersetup{colorlinks=true, linkcolor=blue!60!black, urlcolor=blue!60!black}

\titleformat{\section}[block]{\large\bfseries}{}{0em}{}[{\titlerule[0.4pt]}]
\titleformat{\subsection}[runin]{\bfseries}{}{0em}{}[.\quad]

\newmdenv[
  backgroundcolor=gray!8,
  linecolor=gray!40,
  linewidth=0.6pt,
  innerleftmargin=10pt,
  innerrightmargin=10pt,
  innertopmargin=8pt,
  innerbottommargin=8pt,
  skipabove=6pt,
  skipbelow=4pt
]{resultbox}

\title{\textbf{GX~339-4 HAWK-I Ks-band Reduction}\\[4pt]
       {\large Project Logbook}}
\author{Eoin Trainor \\ University College Cork}
\date{\today}

\begin{document}
\maketitle
\tableofcontents
\newpage

% ─────────────────────────────────────────────────────────────────────────────
\section*{Overview}
\addcontentsline{toc}{section}{Overview}

This logbook documents the full reduction of HAWK-I Ks-band observations of
the black hole low-mass X-ray binary GX~339-4, from raw ESO archive data
through to differential photometry and lightcurve analysis.  Entries are
ordered chronologically and each covers the method applied, the numerical
results obtained, and an interpretation of those results.

The overall goal is to measure the ellipsoidal light variations of the
donor star, which constrain the binary inclination and mass ratio.  This
requires the system to be in deep quiescence, so a major early finding was
identifying which observing blocks were contaminated by an outburst, and
excluding them before photometric analysis.

% ─────────────────────────────────────────────────────────────────────────────
\section{Stage 1 -- Data Preparation and Calibration}

\subsection{Method}
Raw HAWK-I science, dark, and flat-field frames were downloaded from the
ESO science archive.  A Python pipeline organised the files by type,
filter, and detector integration time (DIT).  Master calibration products
were built in the following order:

\begin{enumerate}
  \item \textbf{Master dark} -- sigma-clipped median of dark frames matched
        to the science DIT (10\,s, NDIT\,=\,9).  Bad pixels identified as
        outliers at $>3\sigma$ were flagged in a bad-pixel mask (BPM).
  \item \textbf{Master flat} -- median of Ks twilight flats at DIT\,=\,1.676\,s,
        normalised to unity.  Pixels outside $0.5$--$1.5$ times the median
        were added to the BPM.
  \item \textbf{Science calibration} -- each science frame was
        dark-subtracted, flat-fielded, and bad-pixel corrected.  Only
        chip\,1 (bottom-left detector, where GX~339-4 falls) was processed.
\end{enumerate}

Calibrated frames were aligned within each observing block (OB) using
SExtractor source catalogues and an affine transformation, then grouped
into per-OB stacks.  This produced one co-added image per OB for the
initial photometric exploration.

\begin{figure}[htbp]
  \centering
  \includegraphics[width=\textwidth]{C:/Astronomy/GX 339-4/GX 339-4 Output/visualisation/finding_chart_clean.png}
  \caption{HAWK-I Ks-band reference field of GX~339-4, built from the
           inverse-variance co-addition of all 236 science frames.
           \textit{Left:} full chip\,1 ($2048\times2048$\,px,
           $\sim\!3.6'\times3.6'$); the cyan box marks the $\pm90$\,px
           region around the target.  \textit{Right:} $\pm80$\,px
           ($8.5''$) zoom; GX~339-4 is circled in cyan.  The field is
           severely crowded, with several sources within one PSF FWHM
           of the target.  Scale bars: $1'$ (left) and $5''$ (right).}
  \label{fig:finding_chart}
\end{figure}

\subsection{Results}

\begin{resultbox}
13 OB co-adds produced (OBs 1--9: May--Jul 2025; OBs 10--12: Aug--Sep 2025,
later identified as outburst; OB\,13: a second epoch on 23 Jul 2025).
Median seeing across all frames: $\sim0.62''$ (FWHM $\approx$\,5.8\,px at
$0.106''$/px).
\end{resultbox}

\begin{figure}[htbp]
  \centering
  \includegraphics[width=\textwidth]{C:/Astronomy/GX 339-4/GX 339-4 Output/logs/calibration/summary_plots/summary_01_sky_mean_per_ob.png}
  \caption{Mean Ks-band sky background level per OB.  Error bars show the
           frame-to-frame standard deviation within each OB.  Sky levels
           vary between $\sim\!10\,500$\,ADU (OB\,4, 2025-06-19) and
           $\sim\!18\,000$\,ADU (OB\,11, 2025-09-19) across the campaign,
           reflecting lunar phase and seasonal sky brightness changes.}
  \label{fig:sky_mean}
\end{figure}

\begin{figure}[htbp]
  \centering
  \includegraphics[width=\textwidth]{C:/Astronomy/GX 339-4/GX 339-4 Output/logs/calibration/summary_plots/summary_02_sky_std_per_ob.png}
  \caption{Frame-to-frame sky variability (standard deviation of per-frame
           sky medians) within each OB.  Most OBs show stable sky conditions
           ($\sigma_\mathrm{sky} \lesssim 200$\,ADU).  OB\,9 (2025-07-23)
           is a clear outlier at $\sigma \approx 600$\,ADU, indicating
           variable transparency that night; this is also visible in the
           residuals plot (Fig.~\ref{fig:calib_residuals}).}
  \label{fig:sky_std}
\end{figure}

\begin{figure}[htbp]
  \centering
  \includegraphics[width=\textwidth]{C:/Astronomy/GX 339-4/GX 339-4 Output/logs/calibration/summary_plots/summary_03_residuals_per_ob.png}
  \caption{Per-frame post-sky-subtraction residuals for all 12 OBs.  Each dot
           is one science frame; the horizontal bar is the OB mean.  Residuals
           are centred on zero for OBs~1--11, confirming that sky subtraction
           is unbiased.  OB~9 (2025-07-23) shows elevated scatter, likely
           driven by variable sky transparency that night.}
  \label{fig:calib_residuals}
\end{figure}

\subsection{Interpretation}
The calibration pipeline ran without major issues.  The bad-pixel mask was
dominated by hot pixels from the dark, with the flat adding a small number
of low-sensitivity edge pixels.  Having one co-added image per OB was
sufficient for a first look at the source brightness and for assessing the
BAT X-ray state context.

% ─────────────────────────────────────────────────────────────────────────────
\section{Stage 2 -- Initial Circular Aperture Photometry (13 Co-adds)}

\subsection{Method}
A simple circular aperture was placed on the GX~339-4 centroid in each of
the 13 per-OB co-added images.  The centroid was found by flux-weighted
recentering within a $20\times20$\,px search box around the WCS-predicted
position.  A fixed aperture of radius $r = 6$\,px was used with a sky
background annulus of $r_\mathrm{in} = 10$\,px to $r_\mathrm{out} = 16$\,px.
Background was estimated via iterative sigma-clipping of the annulus pixel
distribution ($3\sigma$, 5 iterations): in each pass, pixels deviating
more than $3\sigma$ from the running median were flagged as outliers and
excluded from subsequent iterations.  This sigma-clip procedure acts as
a per-frame \textit{pixel mask}, suppressing two sources of contamination
that would otherwise bias the sky level upward: (i) PSF wings from the
two bright field stars whose halos overlap the outer annulus in several
OBs, and (ii) isolated hot or bad pixels.  The background level $B$ was
taken as the median of the surviving unmasked pixels; $\sigma_B$ was
their standard deviation, propagated into the SNR estimate.  The net
source flux was the aperture sum minus $B$ times the aperture area.
No comparison-star differential correction was applied at this stage.

\begin{figure}[htbp]
  \centering
  \includegraphics[width=0.48\textwidth]{C:/Astronomy/GX 339-4/GX 339-4 Output/photometry/aperture_v1_outputs/aperture_cutouts/gx3394_0001.png}
  \caption{Ks-band cutout of GX~339-4 (OB\,1, 2025-05-17) with the
           circular aperture ($r = 6$\,px, inner ring) and sky annulus
           ($r = 10$--$16$\,px, outer rings) overlaid.  A bright neighbour
           is visible to the upper right; its halo falls within the outer
           annulus, biasing the sky estimate and contaminating the net flux
           measurement.  The proximity of this neighbour means the
           net flux cannot be cleanly attributed to GX~339-4 alone.}
  \label{fig:ap_cutout}
\end{figure}

\begin{figure}[htbp]
  \centering
  \includegraphics[width=0.48\textwidth]{C:/Astronomy/GX 339-4/GX 339-4 Output/photometry/aperture_v1_outputs/annulus_histograms/annulus_hist_0001.png}
  \hfill
  \includegraphics[width=0.48\textwidth]{C:/Astronomy/GX 339-4/GX 339-4 Output/photometry/aperture_v1_outputs/annulus_histograms/annulus_hist_0011.png}
  \caption{Annulus pixel-value distributions for OB\,1 (left, quiescent,
           2025-05-17) and OB\,11 (right, outburst, 2025-08-14).
           Grey bars show all raw annulus pixels; orange bars show the
           sigma-clipped (retained) subset.  The dashed vertical line is
           the iterative clip threshold; the solid line marks the final
           background median $B$; dotted lines show $B \pm \sigma_B$.
           In OB\,1 the annulus is uncontaminated: the raw and retained
           distributions are nearly identical and $B$ is well-defined.
           In OB\,11 a low-value tail from PSF wings of the saturating
           GX~339-4 source extends into the inner annulus edge; the
           sigma-clip masks these pixels, pulling $B$ toward the true
           sky level and preventing an overestimated background.}
  \label{fig:annulus_hist}
\end{figure}

\subsection{Results}

\begin{resultbox}
\begin{tabular}{lrr}
  \toprule
  Date & Net flux (ADU) & SNR \\
  \midrule
  2025-05-17 (OB\,1) & 27\,688  &  3.8 \\
  2025-06-02 (OB\,2) & 29\,894  &  5.5 \\
  2025-06-04 (OB\,3) & 18\,216  &  3.3 \\
  2025-06-19 (OB\,4) & 33\,723  &  7.7 \\
  2025-07-03 (OB\,5) & 35\,234  &  5.2 \\
  2025-07-14 (OB\,6) & 39\,286  &  5.9 \\
  2025-07-18 (OB\,7) & 51\,766  &  6.8 \\
  2025-07-20 (OB\,8) & 41\,721  & 10.7 \\
  2025-07-23 (OB\,9) & 42\,994  & 10.1 \\
  2025-07-23 (OB\,9b)& 56\,005  &  8.4 \\
  2025-08-14 (OB\,10)& 972\,560 & 141.5 \\
  2025-08-29 (OB\,11)&1\,228\,689& 115.6 \\
  2025-09-19 (OB\,12)& 789\,364 &  97.4 \\
  \bottomrule
\end{tabular}
\end{resultbox}

\begin{figure}[htbp]
  \centering
  \includegraphics[width=\textwidth]{C:/Astronomy/GX 339-4/GX 339-4 Output/photometry/comparison_stars/Comp Stars Residuals/GX 339-4 Relative to Star 3.png}
  \caption{Differential Ks-band magnitude of GX~339-4 relative to comparison
           Star~3 ($K_s \approx 14.5$\,mag) across all 13 per-OB co-adds.
           OBs~1--9 (May--Jul 2025) cluster near $\Delta m \approx 3.3$--$3.6$,
           placing GX~339-4 at $K_s \approx 17$--$18$\,mag.  OBs~10--12
           (Aug--Sep 2025) drop to $\Delta m \approx -0.5$--$+0.1$, indicating
           GX~339-4 has brightened to near the comparison star flux during
           outburst.}
  \label{fig:gx_vs_star3}
\end{figure}

\begin{figure}[htbp]
  \centering
  \includegraphics[width=0.48\textwidth]{C:/Astronomy/GX 339-4/GX 339-4 Output/photometry/aperture_v1_outputs/images/Figure_3.png}
  \hfill
  \includegraphics[width=0.48\textwidth]{C:/Astronomy/GX 339-4/GX 339-4 Output/photometry/aperture_v1_outputs/images/Figure_11.png}
  \caption{Ks-band cutouts with aperture ($r = 6$\,px, inner ring) and sky
           annulus ($r = 10$--$16$\,px, outer rings) overlaid, illustrating
           two distinct contamination regimes.  \textit{Left}: OB\,3
           (2025-06-04, quiescent).  Two bright field stars are clearly
           resolved within the annulus region -- upper-right and lower-right
           of centre -- whose elevated pixels would bias $B$ upward without
           sigma-clip masking.  \textit{Right}: OB\,11 (2025-08-14,
           outburst).  GX~339-4 is saturating and its extended PSF wings
           spill into the inner annulus; the sigma-clip masks the resulting
           high-flux tail, preventing the sky level from being
           overestimated.  In both cases the final background $B$ is
           estimated from the unmasked (orange) pixel subset shown in
           Figure~\ref{fig:annulus_hist}.}
  \label{fig:annulus_cutouts}
\end{figure}

\begin{figure}[htbp]
  \centering
  \includegraphics[width=0.55\textwidth]{C:/Astronomy/GX 339-4/GX 339-4 Output/photometry/aperture_v2_outputs/circular_aperture_outputs/GX 339-4 (Annulus Overlay).png}
  \caption{GX~339-4 aperture and sky annulus with the pixel mask overlaid.
           The yellow circle marks the source aperture; the two cyan circles
           bound the sky annulus ($r = 10$--$16$\,px).  The green region
           shows the \textit{unmasked} (retained) annulus pixels that
           contribute to the sky estimate.  The asymmetric exclusion of
           pixels in the lower portion of the annulus is caused by the
           halo of the bright neighbour (upper left) bleeding into the
           annulus; the sigma-clip procedure identifies these elevated
           pixels and removes them, preventing the sky level from being
           biased high.}
  \label{fig:annulus_overlay}
\end{figure}

\begin{figure}[htbp]
  \centering
  \includegraphics[width=\textwidth]{C:/Astronomy/GX 339-4/GX 339-4 Output/photometry/lightcurve_13ob_aperture.png}
  \caption{GX~339-4 net Ks-band flux vs MJD for all 13 per-OB co-adds
           ($r = 6$\,px aperture, sigma-clipped sky background).
           OBs~1--10 (blue circles, May--Jul 2025) show a quiescent
           baseline of $\sim$20\,000--60\,000\,ADU with a gradual
           upward trend.  OBs~11--13 (red stars, Aug--Sep 2025) jump
           by a factor of $\sim$20--30 to $>900\,000$\,ADU, marking
           the onset of a hard X-ray outburst confirmed by Swift BAT
           (Section~\ref{sec:bat}).  Error bars are $1\sigma$
           (flux/SNR).}
  \label{fig:lc_13ob}
\end{figure}

\subsection{Interpretation}
The most striking feature is a factor of $\sim$20--40 increase in net flux
between OBs~9 and~10 (July to August 2025), with OBs~10--12 showing
fluxes an order of magnitude above the earlier observations.  This
immediately flagged an X-ray outburst in progress during those later OBs.
The low SNR values for OBs~1--9 (3--11) reflect both the faintness of
GX~339-4 in quiescence ($K_s \approx 17$\,mag) and the contamination from
blended field stars within the fixed aperture -- an issue that motivated
the development of a more rigorous aperture approach in Stage~5.
The high SNR for OBs~10--12 is dominated by the bright accretion
disc during outburst, not the donor.

% ─────────────────────────────────────────────────────────────────────────────
\section{Stage 3 -- Swift BAT X-ray Analysis}
\label{sec:bat}

\subsection{Method}
The Swift Burst Alert Telescope (BAT) daily-averaged hard X-ray lightcurve
(15--50\,keV) for GX~339-4 was downloaded from the BAT transient monitor,
covering 2005 to 2026.  Two analyses were carried out:

\begin{enumerate}
  \item \textbf{Overlay of ESO OB windows on the BAT lightcurve} -- Each OB
        start and end time was converted to MJD and plotted as a shaded
        region on the BAT light curve.  This gave immediate visual
        confirmation of which OBs fell in active versus quiescent X-ray
        states.

  \item \textbf{Statistical outburst classification} -- The full BAT
        history was used to define a quiescent baseline via
        high-side-clipping ($1\sigma$ iterative), yielding a quiescent
        rate $\mu_q$ and scatter $\sigma_q$.  Each OB was assigned a
        z-score relative to this baseline, and a coherence-weighted
        neighbour algorithm (requiring $\geq 3\sigma$ points with
        $\geq 50\%$ elevated neighbours within a 2-day window) classified
        each OB as \textit{Quiescent}, \textit{Elevated}, or
        \textit{Outburst}.
\end{enumerate}

\subsection{Results}

\begin{resultbox}
OBs 1--9 (2025 May--Jul): classified as \textbf{Quiescent} or
\textbf{Elevated} -- BAT rates consistent with the low-state baseline.\\[4pt]
OBs 10--12 (2025 Aug--Sep): classified as \textbf{Outburst} -- z-scores of
$10\sigma$--$30\sigma$ above quiescence, coinciding with the factor
$\sim$30 flux increase seen in the NIR aperture photometry.\\[4pt]
Quiescent BAT rate baseline: $\mu_q \approx 0$\,cts\,cm$^{-2}$\,s$^{-1}$
(system below BAT sensitivity in true quiescence).
\end{resultbox}

\begin{figure}[htbp]
  \centering
  \includegraphics[width=\textwidth]{C:/Astronomy/GX 339-4/GX 339-4 Output/xray_analysis/outburst/swift_bat_with_eso_observing_blocks.png}
  \caption{Swift BAT 15--50\,keV daily lightcurve of GX~339-4 (red points)
           over the 2025 campaign window, with ESO HAWK-I observing block
           windows marked as dashed vertical lines (OB numbers labelled at
           the top).  The $\pm 1\sigma$ and $\pm 3\sigma$ quiescence bands
           are shown as dotted lines.  OB\,5 (MJD\,$\approx 60855$) coincides
           with a prominent X-ray flare; OBs~10--12 fall after the outburst
           onset near MJD\,$\approx 60900$.}
  \label{fig:bat_obs}
\end{figure}

\begin{figure}[htbp]
  \centering
  \includegraphics[width=\textwidth]{C:/Astronomy/GX 339-4/GX 339-4 Output/xray_analysis/outburst/analysis_outputs/swift_bat_with_eso_observing_blocks_quiescence_sigma.png}
  \caption{As Fig.~\ref{fig:bat_obs}, with the quiescence baseline
           ($\mu_q$, blue solid line) and $+3\sigma$ outburst threshold
           (blue dashed) derived from high-side sigma-clipping of the full
           BAT history.  Points flagged as $\geq 3\sigma$ outburst
           detections are highlighted in orange.  OBs~10--12 land firmly
           above the outburst threshold.}
  \label{fig:bat_sigma}
\end{figure}

\subsection{Interpretation}
The BAT analysis conclusively showed that the system underwent a hard X-ray
outburst beginning around 2025 August, coinciding exactly with the dramatic
NIR brightening seen in Stage~2.  During outburst the accretion disc
dominates the NIR flux and drowns any ellipsoidal signal from the donor
star.  OBs~10--12 were therefore excluded from all subsequent photometric
analysis.  OBs~1--9 (May--July 2025) were retained as the quiescent
dataset, though the elevated BAT classification for some of these OBs
suggested the system may not have been in deep quiescence even then, which
partly explains the failure to detect a clean ellipsoidal signal later.

% ─────────────────────────────────────────────────────────────────────────────
\section{Stage 4 -- Orbital Phase Coverage Analysis}

\subsection{Method}
The mid-exposure MJD of each OB co-add was computed from the FITS headers
and converted to orbital phase using the Heida et al.\ (2017) ephemeris:
\begin{equation}
  \phi = \left(\frac{\mathrm{MJD} - T_0}{P}\right) \bmod 1,
  \qquad T_0 = 57529.397\;\mathrm{MJD},\quad P = 1.7587\;\mathrm{d}.
\end{equation}
Here $\phi = 0$ corresponds to inferior conjunction of the donor star
(donor closest to the observer, ellipsoidal lightcurve minimum), and
$\phi = 0.25$ and $\phi = 0.75$ correspond to quadrature (ellipsoidal
maxima).

A phase coverage histogram was generated binning the observations into 100
equal intervals over the full $\phi = 0$--$1$ range.

\subsection{Results}

\begin{resultbox}
\begin{tabular}{lrl}
  \toprule
  OB & Date & Phase $\phi$ \\
  \midrule
  OB\,1  & 2025-05-17 & 0.630 \\
  OB\,2  & 2025-06-02 & 0.740 \\
  OB\,3  & 2025-06-04 & 0.847 \\
  OB\,4  & 2025-06-19 & 0.295 \\
  OB\,5  & 2025-07-03 & 0.284 \\
  OB\,6  & 2025-07-14 & 0.579 \\
  OB\,7  & 2025-07-18 & 0.824 \\
  OB\,8  & 2025-07-20 & 0.924 \\
  OB\,9  & 2025-07-23 & 0.597--0.605 \\
  \bottomrule
\end{tabular}\\[6pt]
Phase range sampled (quiescent OBs 1--9): $\phi \approx 0.28$--$0.92$.\\
Gap: inferior conjunction region ($\phi \lesssim 0.28$) not observed.
\end{resultbox}

\begin{figure}[htbp]
  \centering
  \includegraphics[width=\textwidth]{C:/Astronomy/GX 339-4/GX 339-4 Output/xray_analysis/orbital_phase/phase_coverage_clean.png}
  \caption{Orbital phase coverage of the 234 good frames across 9 OBs,
           binned into 20 equal phase intervals ($\Delta\phi = 0.05$).
           Blue bars show frames in covered bins (13--48 frames per bin);
           the shaded red region ($\phi < 0.28$) marks the unobserved gap
           around inferior conjunction.  Blue dotted lines indicate
           conjunction phases ($\phi = 0, 0.5$) and red dotted lines
           indicate quadrature ($\phi = 0.25, 0.75$), where the
           ellipsoidal lightcurve reaches its minima and maxima
           respectively.}
  \label{fig:phase_coverage}
\end{figure}

\subsection{Interpretation}
While 9 OBs span a wide range of phases, the region around inferior
conjunction ($\phi \sim 0$--$0.28$) is entirely absent from the dataset.
This is a significant gap: the ellipsoidal minima at $\phi = 0$ and
$\phi = 0.5$ are needed to anchor the lightcurve shape and amplitude, and
without the $\phi \sim 0$ minimum it is impossible to fully constrain the
ellipsoidal model.  Phase coverage is clustered in the quadrature and
post-conjunction ranges.  Future observations should specifically target
OB windows that land near $\phi \approx 0.0$--$0.2$.

% ─────────────────────────────────────────────────────────────────────────────
\section{Stage 5 -- Refined Aperture Photometry with Comparison Stars}

\subsection{Method}
A more rigorous aperture photometry pipeline was developed for
single-frame analysis, motivated by the low SNR and field-crowding problems
seen in Stage~2.  The pipeline proceeded as:

\begin{enumerate}
  \item \textbf{PSF measurement} -- Five comparison stars spanning
        $K_s \approx 14$--$17$\,mag were selected.  Each was fitted with
        a 2D Gaussian; the median FWHM across all five defined the
        reference PSF scale.
  \item \textbf{Annulus geometry} -- Sky annulus fixed at $3.0$--$5.0$
        times the median FWHM ($22.7$--$37.8$\,px), ensuring consistent
        background sampling across all stars.
  \item \textbf{Background pixel masking (truncated-core method)} --
        Rather than iterative sigma-clipping, sky background was estimated
        from the \textit{core} of the annulus pixel distribution.
        The mode of the distribution was located via MAD-based robust
        statistics, and a symmetric window
        $[B_\mathrm{mode} - \delta,\; B_\mathrm{mode} + \delta]$ was
        placed around it.  Any annulus pixel falling outside this window
        was masked and excluded from the background estimate.  This
        rejects PSF halos from neighbouring stars, bad pixels, and
        scattered light without relying on iterative convergence.
        The background $B$ was the mean of surviving core pixels; the
        SNR calculation used the core pixel count $N_\mathrm{core}$
        (rather than the full annulus area) to avoid overstating the
        effective sample size.  For OB\,1, this reduced 3116 raw annulus
        pixels to 2309 core pixels (74\%), collapsing the pixel spread
        from a raw standard deviation of 3318\,ADU to a core MAD of
        6.4\,ADU -- a factor $\sim$500 improvement in background
        precision.
  \item \textbf{Comparison star ranking} -- Stars were ranked by maximum
        SNR$_3$ (SNR using core-pixel background estimation), with
        penalties for saturation, poor Gaussian fit, and PSF axis ratio.
  \item \textbf{GX aperture selection} -- The optimal aperture for GX~339-4
        was found by a stability criterion: the smallest radius at which
        $|\Delta F(r_i \to r_{i+1})| \leq 2\sqrt{\sigma_i^2 + \sigma_{i+1}^2}$
        for two consecutive steps.  If no plateau was found, the
        peak-SNR aperture was used as a fallback.
\end{enumerate}

\begin{figure}[htbp]
  \centering
  \includegraphics[width=0.48\textwidth]{C:/Astronomy/GX 339-4/GX 339-4 Output/photometry/aperture_v2_outputs/circular_aperture_outputs/Annulus Histogram.png}
  \hfill
  \includegraphics[width=0.48\textwidth]{C:/Astronomy/GX 339-4/GX 339-4 Output/photometry/aperture_v2_outputs/aperture_outputs_final/GX 339-4 Annulus Overlay.png}
  \caption{\textit{Left}: Annulus pixel-value histogram for OB\,1
           (2025-05-17) using the v2 pipeline.  The dominant spike at
           $\approx 14\,200$\,ADU is the true sky background; a small
           cluster near 0\,ADU arises from bad-pixel regions that fall
           within the large annulus.  The solid blue line marks the core
           mean $B$; dotted lines bound the tight core window used for
           masking.  All pixels outside this window are excluded before
           $B$ is computed.
           \textit{Right}: GX~339-4 field with photometric aperture
           (yellow circle, $r = 5.59$\,px) and sky annulus (cyan rings,
           $r = 22.7$--$37.8$\,px) overlaid.  The dark green shading
           marks the annulus pixels that survived the truncated-core mask
           and were used for background estimation; unshaded areas within
           the annulus rings (notably toward the lower-right) were masked
           out.  The bright field star visible at upper-left lies outside
           the annulus.}
  \label{fig:core_mask}
\end{figure}

\subsection{Results}

\begin{resultbox}
Median FWHM from 5 comparison stars: 7.55\,px ($0.80''$).\\[3pt]
\textbf{Comparison star ranking:}\\
\begin{tabular}{lrrrr}
  \toprule
  Star & FWHM (px) & Max SNR$_3$ & Opt $r$ (px) & Score \\
  \midrule
  Star 3 & 7.42 & 9916 & 5.00 & 16458 \\
  Star 2 & 7.63 & 8783 & 5.00 & 14594 \\
  Star 1 & 7.57 & 7437 & 5.00 & 12330 \\
  Star 4 & 7.55 & 7287 & 5.00 & 12097 \\
  Star 5 & 2.13 & 1008 & 16.0 & 309   \\
  \bottomrule
\end{tabular}\\[6pt]
\textbf{GX 339-4 measurement (OB\,1 co-add):}\\
Chosen aperture: $r = 5.59$\,px (max-SNR fallback -- no plateau found).\\
Net flux: 26\,557\,ADU.  SNR$_3 = 413$.  Background: 14\,215\,ADU\,px$^{-1}$.
\end{resultbox}

\begin{figure}[htbp]
  \centering
  \includegraphics[width=\textwidth]{C:/Astronomy/GX 339-4/GX 339-4 Output/photometry/aperture_v2_outputs/circular_aperture_outputs/GX 339-4 Analysis.png}
  \caption{Four-panel aperture diagnostic for GX~339-4 (OB\,1,
           2025-05-17).
           \textit{Top left}: field cutout with chosen aperture
           ($r = 5.59$\,px, yellow) and sky annulus
           ($r = 22.7$--$37.8$\,px, cyan).
           \textit{Top right}: radial mean (blue) and median (orange) ADU
           profile.  Both flatten to the sky level
           ($\approx 14\,200$\,ADU) by $\sim$20\,px, confirming the
           annulus inner edge begins well beyond the PSF core; the
           mean exceeds the median only at small radii where the PSF
           peak dominates.
           \textit{Bottom left}: curve of growth (background-subtracted
           cumulative flux vs aperture radius).  The continuous rise
           with no plateau confirms that blended neighbour flux
           contaminates the aperture at all radii.
           \textit{Bottom right}: GX~339-4 SNR$_3$ vs aperture radius;
           the peak at $r = 5.59$\,px defines the chosen aperture via
           the max-SNR fallback.}
  \label{fig:gx_analysis}
\end{figure}

\begin{figure}[htbp]
  \centering
  \includegraphics[width=0.8\textwidth]{C:/Astronomy/GX 339-4/GX 339-4 Output/photometry/aperture_v2_outputs/aperture_outputs_final/Stability Test.png}
  \caption{GX~339-4 aperture stability test (OB\,1).  Blue points show
           $|\Delta F|$ between adjacent aperture radii; the orange
           dashed line shows $2\times$ the combined photometric
           uncertainty (the stability threshold).  A plateau would be
           flagged at the first radius where two consecutive blue points
           fall below the orange line; no such crossing occurs across
           the full tested range ($r = 2.6$--$18.7$\,px).  This failure
           directly demonstrates that blended neighbour flux adds
           continuously with aperture size, making the max-SNR fallback
           ($r = 5.59$\,px) the only viable aperture choice.}
  \label{fig:stability_test}
\end{figure}

\subsection{Interpretation}
The refined pipeline achieved SNR\,$\approx 413$ on the OB\,1 co-add for
GX~339-4, a substantial improvement over the $\sim$4 obtained in Stage~2
for the same OB.  The main gains came from the much larger annulus
(22.7--37.8\,px versus 10--16\,px), which avoids contamination from the
halo of nearby stars, and from the tighter background estimation using only
the core of the annulus pixel distribution.

Star~5 was an outlier: its measured FWHM of 2.1\,px is unrealistically
small (likely a failed Gaussian fit), and its optimal aperture radius of
16\,px is much larger than for the other stars, suggesting it is blended or
saturated.  It was penalised in the ranking and not used.

The failure of the stability criterion to find a plateau for GX~339-4 is
significant: the flux continues rising with aperture radius because of
blended light from a nearby field star.  This means no aperture radius
provides a clean, uncontaminated measurement of GX~339-4's flux alone ---
any net flux value carries an irreducible systematic contribution from
the blended neighbour, limiting the reliability of this photometry.

\begin{figure}[htbp]
  \centering
  \includegraphics[width=0.7\textwidth]{C:/Astronomy/GX 339-4/GX 339-4 Output/photometry/aperture_v2_outputs/aperture_outputs_final/Flux vs Aperture Radius.png}
  \caption{Net flux of GX~339-4 as a function of circular aperture radius
           (OB\,1 co-add).  Unlike an isolated point source, the flux rises
           continuously with no plateau, because contaminating light from a
           blended neighbour enters the aperture at all radii.  The chosen
           aperture ($r = 5.59$\,px, dashed line) corresponds to the
           peak-SNR fallback; no stability criterion plateau was found.}
  \label{fig:flux_vs_radius}
\end{figure}

\begin{figure}[htbp]
  \centering
  \includegraphics[width=\textwidth]{C:/Astronomy/GX 339-4/GX 339-4 Output/photometry/curve_of_growth.png}
  \caption{\textit{Left:} Normalised cumulative curve of growth for four
           reference stars (blue shades, $K_s = 9.5$--$11.2$\,mag) and
           GX~339-4 (red), measured on the deep reference image.  All
           curves are normalised to 1.0 at $r = 15$\,px.  The reference
           stars converge smoothly to unity by $r \approx 10$\,px,
           as expected for isolated point sources.  GX~339-4 rises
           continuously and exceeds 1.0, confirming that flux from the
           blended neighbour enters the aperture at all radii.
           \textit{Right:} SNR$_3$ vs aperture radius for the four
           comparison stars from the refined single-frame analysis.
           All four peak near $r = 5$\,px (marked), consistent with
           the PSF FWHM of $\sim$7.5\,px; choosing the peak-SNR
           aperture is well motivated for isolated sources.}
  \label{fig:cog}
\end{figure}

% ─────────────────────────────────────────────────────────────────────────────
\section{Stage 6 -- ZOGY Proper Difference Imaging}

\subsection{Method}
ZOGY proper image subtraction (Zackay, Ofek \& Gal-Yam 2016) was applied
to all 9 quiescent OBs (OBs~1--9, 236 frames total).  The algorithm
proceeds in Fourier space so that the difference image $\hat{D}$ has
spatially uniform noise:

\begin{equation}
  \hat{D} = \frac{F_r\,\hat{P}_r\,\hat{N} - F_n\,\hat{P}_n\,\hat{R}}
                 {\sqrt{F_r^2\,\sigma_n^2\,|\hat{P}_r|^2 +
                        F_n^2\,\sigma_r^2\,|\hat{P}_n|^2}},
\end{equation}

where $R$ and $N$ are the reference and science images, $P_r$, $P_n$ their
PSFs, $\sigma_r$, $\sigma_n$ background RMS values, and $F_r$, $F_n$
photometric flux normalisations.  The implementation steps were:

\begin{enumerate}
  \item \textbf{Reference construction} -- A deep reference image $R$ was
        built by inverse-variance-weighted co-addition of all 9 OBs in
        Fourier space, matching each frame's PSF to a common kernel.

  \item \textbf{PSF modelling} -- Per-frame PSFs were modelled as Moffat
        profiles fitted to isolated bright stars.  Moffat profiles were
        adopted over Gaussians because they better reproduce the extended
        wings of seeing-broadened stellar PSFs, reducing the residual
        power in $D$ at large separations from subtracted sources.
        Flux zero-points $F_n$ were determined by matching instrumental
        magnitudes to 2MASS $K_s$ catalogue values.

  \item \textbf{Difference image and $S$ statistic} -- Each of the 236
        science frames was subtracted against the reference to produce $D$.
        The matched-filter detection statistic $S$ was computed by
        convolving $\hat{D}$ with the science PSF in Fourier space.
        $S_\mathrm{target}$ at the GX~339-4 centroid gives an SNR-optimal
        estimate of the differential flux per frame.

  \item \textbf{Reference star check} -- $S$ was also extracted at eight
        non-variable reference star positions as a systematics diagnostic.
\end{enumerate}

\begin{figure}[htbp]
  \centering
  \includegraphics[width=\textwidth]{C:/Astronomy/GX 339-4/GX 339-4 Output/difference/report_figures/fig1_reference_image.png}
  \caption{ZOGY reference image $R$ built from the inverse-variance-weighted
           co-addition of all 236 science frames (9 OBs).  \textit{Left:}
           full HAWK-I chip\,1 field ($2048 \times 2048$\,px,
           $\sim 3.6' \times 3.6'$); the target region is indicated by the
           red box.  \textit{Right:} 80\,px ($8.5''$) zoom on GX~339-4
           (circled), showing the crowded field environment.}
  \label{fig:reference_image}
\end{figure}

\begin{figure}[htbp]
  \centering
  \includegraphics[width=\textwidth]{C:/Astronomy/GX 339-4/GX 339-4 Output/difference/report_figures/fig2_before_after.png}
  \caption{ZOGY before/after comparison for a single science frame
           (OB\,1, frame 1).  \textit{(a)} Science frame $N$: field stars
           and GX~339-4 visible in the crowded field.
           \textit{(b)} Difference image $D$: field stars subtract cleanly
           to noise, leaving only sources that changed flux relative to the
           reference.  GX~339-4 (crosshair) shows a residual consistent
           with its differential flux at this epoch.
           \textit{(c)} Reference image $R$: the deep co-add used for
           subtraction.  All panels show an 80\,px zoom around GX~339-4.}
  \label{fig:before_after}
\end{figure}

\subsection{Results}

\begin{resultbox}
236 frames processed; 2 rejected by per-OB $3\sigma$ clipping
(1 each in OBs~2 and 8), leaving \textbf{234 good frames}.\\[3pt]
Median PSF FWHM (Moffat): 6.00\,px = $0.64''$ (range $\approx 0.37''$--$0.94''$) across all OBs.\\[3pt]
\begin{tabular}{lcc}
  \toprule
  OB & Frames (good) & Mean FWHM ($''$) \\
  \midrule
  OB\,1 & 26/26 & 0.85 \\
  OB\,2 & 25/26 & 0.82 \\
  OB\,3 & 26/26 & 0.55 \\
  OB\,4 & 26/26 & 0.76 \\
  OB\,5 & 26/26 & 0.70 \\
  OB\,6 & 26/26 & 0.75 \\
  OB\,7 & 26/26 & 0.61 \\
  OB\,8 & 25/26 & 0.83 \\
  OB\,9 & 28/28 & 0.55 \\
  \bottomrule
\end{tabular}\\[6pt]
Raw $S_\mathrm{target}$ scatter: $\sigma \approx 8 \times 10^8$ matched-filter units.\\
OB-to-OB offset dominates over intra-OB variability by $\sim$factor of 10.
\end{resultbox}

\begin{figure}[htbp]
  \centering
  \includegraphics[width=\textwidth]{C:/Astronomy/GX 339-4/GX 339-4 Output/logs/zogy/quality/q02_fwhm_timeline.png}
  \caption{Moffat-fitted PSF FWHM (pixels) for each of the 234 science
           frames, ordered by frame index.  Vertical grey lines separate
           OBs.  The median FWHM is 6.00\,px ($0.64''$); individual frames
           span approximately 3.5--8.9\,px ($0.37''$--$0.94''$).  The
           large seeing variation between and within OBs drives the
           PSF-matching residuals that dominate the ZOGY noise floor.}
  \label{fig:fwhm_timeline}
\end{figure}

\begin{figure}[htbp]
  \centering
  \includegraphics[width=\textwidth]{C:/Astronomy/GX 339-4/GX 339-4 Output/logs/zogy/quality/q03_zeropoint_timeline.png}
  \caption{Photometric flux zero-point $F_n$ per frame, derived by
           matching instrumental aperture fluxes to 2MASS $K_s$
           magnitudes.  The median is $1.17\times10^{10}$\,ADU per
           2MASS linear flux unit.  The factor $\sim$3--4 scatter within
           some OBs reflects variable atmospheric transparency and
           changing PSF encircled energy, and contributes to the
           differential-flux uncertainty in the $S$ statistic.}
  \label{fig:zp_timeline}
\end{figure}

\begin{figure}[htbp]
  \centering
  \includegraphics[width=\textwidth]{C:/Astronomy/GX 339-4/GX 339-4 Output/logs/zogy/quality/q05_diff_pixel_stats.png}
  \caption{Pixel statistics of the ZOGY difference images across all 234
           frames.  \textit{Left:} distribution of $\sigma(D)$, the
           $3\sigma$-clipped background standard deviation of each $D$
           image.  The ZOGY formalism predicts $\sigma(D) = 1$ for a
           correctly normalised difference image; the observed median of
           0.738 (up from 0.726 with the previous Gaussian PSF model)
           reflects the improvement from Moffat wing fitting, but remains
           below the ideal value of 1.0.  The remaining deficit
           ($\approx 0.26$) is attributed to spatially varying PSF across
           the HAWK-I mosaic, sub-pixel centring errors, and
           flat-fielding-induced correlated noise, none of which are
           corrected by the PSF profile choice alone.
           \textit{Right:} distribution of $\mathrm{mean}(D)$, centred
           on zero (median $= 0.036$), confirming no systematic bias in
           the sky subtraction.}
  \label{fig:diff_pixel_stats}
\end{figure}

\begin{figure}[htbp]
  \centering
  \includegraphics[width=\textwidth]{C:/Astronomy/GX 339-4/GX 339-4 Output/difference/report_figures/fig3_sample_diff.png}
  \caption{Sample ZOGY difference image $D$ for OB\,1, frame 1.
           \textit{Left:} full chip ($2048 \times 2048$\,px); the stretch
           is $\pm 5\sigma_D$.  Most of the field subtracts cleanly to
           noise; a handful of residuals from poorly-matched PSFs are
           visible as blue/red dipoles.  \textit{Right:} 80\,px zoom
           on GX~339-4, showing the faint differential residual.}
  \label{fig:sample_diff}
\end{figure}

\begin{figure}[htbp]
  \centering
  \includegraphics[width=\textwidth]{C:/Astronomy/GX 339-4/GX 339-4 Output/logs/zogy/quality/q04_reference_psf.png}
  \caption{Reference PSF $P_R$ of the co-added reference image (236
           frames).  \textit{Left:} 2D image of the effective PSF kernel.
           \textit{Centre:} azimuthally averaged radial profile on a
           logarithmic scale; the FWHM/2 half-power radius is marked (red
           dashed).  \textit{Right:} row and column cross-sections through
           the PSF centre.  The reference PSF has FWHM $\approx 5.60$\,px
           ($0.59''$), consistent with the median Moffat-fitted seeing
           of the campaign.}
  \label{fig:ref_psf}
\end{figure}

\subsection{Interpretation}
ZOGY successfully removed the blended field-star contamination that
defeated aperture photometry, since field stars cancel in $D$ (they are
present identically in $R$ and $N$).  The per-frame $S_\mathrm{target}$
values are physically meaningful estimates of the differential flux of
GX~339-4 relative to the reference epoch.  Switching from Gaussian to
Moffat PSF profiles improved the median $\sigma(D)$ from 0.726 to 0.738,
a small but real gain from better wing modelling; the residual deficit
from the ideal value of 1.0 is dominated by spatially varying PSF,
sub-pixel centring errors, and correlated noise from the reduction
pipeline.

However, the large OB-to-OB variation in $S_\mathrm{target}$ indicates
that the mean NIR flux of the system changed significantly between OBs,
consistent with a semi-active accretion state rather than true quiescence.
The reference stars showed comparable OB-to-OB scatter, confirming that
PSF-matching residuals (driven by variable seeing between OBs) contribute
to the systematics.  This OB-level variability is the dominant noise
source and sets the floor below which the ellipsoidal signal
($\Delta K_s \sim 0.1$--$0.2$\,mag in quiescence) cannot be recovered
without detrending.

% ─────────────────────────────────────────────────────────────────────────────
\section{Stage 7 -- Phase-Folded Lightcurve}

\subsection{Method}
The 234-frame $S_\mathrm{target}$ time series was processed into a
phase-folded lightcurve via the following steps:

\begin{enumerate}
  \item \textbf{Per-OB sigma clipping} -- Frames deviating by more than
        $3\sigma$ from their OB median $S_\mathrm{target}$ were rejected
        (removes cosmic-ray hits and tracking failures within an OB).

  \item \textbf{Per-OB detrending} -- The mean $S_\mathrm{target}$ of each
        OB was subtracted.  This removes the OB-to-OB flux offsets
        (accretion variability, photometric zero-point drift) while
        preserving intra-OB variability.  Each OB spans $\sim$52\,min
        $\ll P = 1.7587$\,d, so the ellipsoidal slope within a single OB
        is negligible.

  \item \textbf{Normalisation} -- The detrended series was divided by its
        global standard deviation to give dimensionless units.

  \item \textbf{Phase folding} -- Orbital phases were computed using the
        Heida+2017 ephemeris and the data binned into 20 equal phase
        intervals (minimum 3 points per bin required for a bin to be plotted).
\end{enumerate}

\subsection{Results}

\begin{resultbox}
234 frames, 9 OBs, 66.8-day baseline ($\approx 38$ orbital periods).\\[3pt]
Phase coverage: $\phi \approx 0.28$--$0.94$; inferior conjunction region
($\phi \lesssim 0.28$) \textbf{not sampled}.\\[3pt]
All binned means consistent with zero within uncertainties.\\
\textbf{No significant ellipsoidal modulation detected.}\\[3pt]
Frame-to-frame scatter: $\sigma(S_\mathrm{detrended}) = 1.0$ (by construction);
individual frames span $-3.5$ to $+3.2$ normalised units.
\end{resultbox}

\begin{figure}[htbp]
  \centering
  \includegraphics[width=\textwidth]{C:/Astronomy/GX 339-4/GX 339-4 Output/difference/lightcurve/lc01_raw_vs_mjd.png}
  \caption{Raw ZOGY matched-filter statistic $S_\mathrm{target}$ vs MJD for
           all 234 good frames, coloured by OB.  Dashed horizontal lines
           show the per-OB mean.  The large OB-to-OB variation (means
           spanning $-4\times10^8$ to $+1.3\times10^9$ units) reflects
           variable accretion activity across the 66.8-day campaign; this
           long-term trend must be removed before phase-folding.}
  \label{fig:lc_raw}
\end{figure}

\begin{figure}[htbp]
  \centering
  \includegraphics[width=\textwidth]{C:/Astronomy/GX 339-4/GX 339-4 Output/difference/lightcurve/lc02_reference_stars.png}
  \caption{ZOGY $S$ statistic for 8 non-variable reference stars
           ($K_s = 9.5$--$11.7$\,mag) vs MJD.  Non-variable sources
           should produce $S \approx 0$ in all frames; the observed
           OB-to-OB scatter (amplitude $\sim 10^9$ units) matches that
           of the GX~339-4 target, confirming the noise is dominated by
           PSF-matching systematics common to all sources in the field,
           not by astrophysical variability in the target itself.}
  \label{fig:ref_stars}
\end{figure}

\begin{figure}[htbp]
  \centering
  \includegraphics[width=\textwidth]{C:/Astronomy/GX 339-4/GX 339-4 Output/difference/lightcurve/lc03_detrended_vs_mjd.png}
  \caption{Detrended differential $S$ statistic vs MJD after subtracting
           the per-OB mean.  The long-term trend has been removed and the
           residuals are centred on zero within each OB.  The remaining
           scatter ($\sigma \approx 1$ after normalisation) represents the
           intra-OB variability from which the phase-folded ellipsoidal
           signal must be extracted.}
  \label{fig:lc_detrended}
\end{figure}

\begin{figure}[htbp]
  \centering
  \includegraphics[width=\textwidth]{C:/Astronomy/GX 339-4/GX 339-4 Output/difference/lightcurve/lc05_phase_folded_paper.png}
  \caption{Phase-folded Ks-band differential lightcurve of GX~339-4
           on the Heida+2017 ephemeris ($P = 1.7587$\,d).  Individual
           frames are coloured by OB; black points with error bars are
           binned means ($\pm$ standard error, 20 bins).  Blue dotted
           lines: inferior/superior conjunction; red dotted lines:
           quadrature.}
  \label{fig:lc}
\end{figure}

\begin{figure}[htbp]
  \centering
  \includegraphics[width=\textwidth]{C:/Astronomy/GX 339-4/GX 339-4 Output/difference/lightcurve/lc06_ob_means_vs_phase.png}
  \caption{\textit{Left:} per-OB mean $S_\mathrm{target}$ vs MJD, showing
           the long-term rising trend from OB\,1 to OB\,9 (a factor
           $\sim$4 increase over 67\,days), indicative of growing accretion
           activity rather than a stable quiescent baseline.
           \textit{Right:} the same nine OB means phase-folded on the
           Heida+2017 ephemeris.  No coherent ellipsoidal pattern is
           visible; the dominant signal is the long-term trend that was
           removed during per-OB detrending.}
  \label{fig:ob_means}
\end{figure}

\subsection{Interpretation}
The non-detection of an ellipsoidal signal is not unexpected given the
evidence that the system was not in deep quiescence during this campaign.
Three compounding factors explain the result:

\begin{enumerate}
  \item \textbf{Accretion variability.}  The large OB-to-OB $S_\mathrm{target}$
        variation ($\sigma \sim 8 \times 10^8$ units) indicates that the
        accretion disc was contributing variable flux at a level that
        exceeds the expected ellipsoidal amplitude.  The detrending removes
        this inter-OB variability but also removes any long-term ellipsoidal
        signal that may be present.

  \item \textbf{Incomplete phase coverage.}  The inferior conjunction region
        ($\phi \lesssim 0.28$) is entirely absent.  A two-minimum ellipsoidal
        curve requires sampling both $\phi \approx 0$ and $\phi \approx 0.5$
        to anchor the minimum flux level.  Without $\phi \sim 0$, the
        amplitude cannot be measured even if the signal is present.

  \item \textbf{PSF systematics.}  Residual PSF-matching errors between OBs
        (driven by seeing changes from $0.42''$ to $1.03''$) create
        correlated noise in $S$ that mimics astrophysical variability and
        increases the scatter within each phase bin.
\end{enumerate}

The pipeline itself is validated: the reference stars show the same noise
level as the target, confirming that the subtraction is working correctly
and the scatter is not from instrumental artefacts specific to the target
position.

% ─────────────────────────────────────────────────────────────────────────────
\section{Summary and Next Steps}

\subsection{What Was Achieved}

A full end-to-end Ks-band reduction pipeline was developed and applied to
9 quiescent HAWK-I observing blocks (234 frames, 66.8-day baseline).  Key
milestones:

\begin{itemize}
  \item Built a complete calibration pipeline (dark, flat, BPM,
        calibration, alignment) for HAWK-I chip\,1 data.
  \item Identified an X-ray outburst in OBs~10--12 via Swift BAT,
        preventing contaminated data from entering the photometric analysis.
  \item Demonstrated that aperture photometry is insufficient for GX~339-4
        due to field crowding, motivating the ZOGY approach.
  \item Implemented ZOGY difference imaging with per-frame PSF modelling,
        2MASS flux calibration, and matched-filter detection statistics.
  \item Produced a phase-folded differential lightcurve using 234
        per-frame measurements across 9 OBs.
\end{itemize}

\subsection{Why No Ellipsoidal Signal Was Detected}

The dominant obstacle is accretion activity: even the ``quiescent'' OBs~1--9
show OB-to-OB flux variation at a level inconsistent with a system in deep
quiescence.  The per-OB detrending necessary to remove this variability
destroys any ellipsoidal signal that is coherent across the full
66.8-day baseline.

\subsection{Required for a Successful Ellipsoidal Detection}

\begin{enumerate}
  \item \textbf{True quiescence} -- Observations during a confirmed X-ray
        quiescent state (BAT rate $< 2\sigma$ above baseline for $\geq 1$
        month before observations), where the donor dominates the NIR flux.
  \item \textbf{Complete phase coverage} -- OBs that sample $\phi \approx
        0.0$--$0.3$ (currently entirely missing).
  \item \textbf{Higher cadence within OBs} -- Longer OBs ($\sim$3--4\,hr)
        would allow intra-OB phase coverage of $\Delta\phi \sim 0.1$,
        making the ellipsoidal slope measurable within a single night.
\end{enumerate}

% ─────────────────────────────────────────────────────────────────────────────
\begin{thebibliography}{9}

\bibitem{Heida2017}
  Heida M. et al., 2017, ApJ, 846, 132.

\bibitem{Markert1973}
  Markert T.~H. et al., 1973, ApJ, 184, L67.

\bibitem{Zackay2016}
  Zackay B., Ofek E.~O., Gal-Yam A., 2016, ApJ, 830, 27.

\end{thebibliography}

\end{document}
